\subsection{Ring theory}

A \emph{ring} is a set $R$ equipped with two binary operations, addition ($+$) and multiplication ($\cdot$), satisfying the following three sets of axioms, called the ring axioms: $R$ is an abelian group under addition, $R$ is a monoid under multiplication, and multiplication is distributive with respect to addition. We further assume the addition and multiplication to be commutative and with unity. The most familiar example of a ring is the ring~$\ZZ$ of rational integers. 

Given a ring $R$, a subset $I$ of $R$ is said to be an \emph{ideal} if $I$ is an additive subgroup of the additive group of $R$ that absorbs multiplication by the elements of $R$. 
Given a rational prime~$p\in \ZZ$, the additive group $p\ZZ$ is an ideal of $\ZZ$.
Any ideal $I$ of $R$ that is not the whole of $R$ is said to be a \emph{proper ideal}, that is, the underlying set of $I$ is a proper subset of the underlying set of $R$. A proper ideal $I$ is called a \emph{prime ideal} if for any $a$ and $b$ in $R$, if $ab$ is in $I$, then at least one of $a$ and $b$ is in $I$.

An \emph{$R$-module} over a ring $R$ is a generalisation of the notion of vector space over a field, wherein scalars are elements of a given ring with identity and an operation of multiplication (on the left and/or on the right), called scalar multiplication,  defined between elements of the ring and elements of the module.

\subsection{Algebraic number theory}
A complex number $\alpha$ is \emph{algebraic} if it is a root of a univariate polynomial with integer coefficients. The defining polynomial of $\alpha$, denoted $f_\alpha$, is the unique (up to multiplication by $\pm 1$) integer polynomial of least degree, whose coefficients have no common factor, that has $\alpha$ as a root.
The \emph{degree} of an algebraic number $\alpha$ is the degree of its minimal polynomial $f_\alpha$.
If $f_\alpha$ is monic then we say that $\alpha$ is an \emph{algebraic integer}.
The sum, the difference, the product and the quotient of two algebraic numbers (except for division by zero) are algebraic numbers; this means that the set of all algebraic numbers is a \emph{field}, commonly denoted by $\Bar{\QQ}$. The sum, the difference, and the product of two algebraic integers is again an algebraic integer; given an algebraic field $K$, the algebraic integers of $K$, form a ring denoted $\OO_K$, called the ring of integers.

A field $K$ is said to be a \emph{field extension}, denoted $K/L$, of a field $L$, if $L$ is a subfield of $K$. Given a field extension $K/L$, the larger field $K$ is an $L$-vector space. The dimension of this vector space is called the \emph{degree} of the extension and is denoted by $[K : L]$.

An \emph{algebraic number field} (or simply number field) $K$ is a finite degree field extension of the field of rational numbers $\QQ$.
Thus $K$ is a field that contains $\QQ$ and has finite dimension when considered as a vector space over $\QQ$. 
It is well-known that each number field $K$ is a simple extension of $\QQ$, i.e., $K$ can be represented as $K = \QQ(\alpha)$, which is generated by the adjunction of a single element $\alpha \in K$, which is said to be the \emph{primitive element}.
 
The Gaussian rationals~$\QQ(i)$  are the first nontrivial example of an algebraic number field, obtained by adjoining $i:=\sqrt{-1}$ to~$\QQ$.  All elements of~$\QQ(i)$ can be written as expressions of the form $a+bi$ with $a,b\in \QQ$; hence $[\QQ(i) : \QQ]=2$. 
Furthermore, $\OO_{\QQ(i)}:=\ZZ[i]$.

An \emph{order}~$\OO$ in a number field $K$ is a free $\ZZ$-submodule of $\OO_K$ of rank $[K:\QQ]$. Since $\OO_K$ is also a free $\ZZ$-module of rank $[K:\QQ]$, it follows from the structure theorem for $\ZZ$-modules that the quotient $\OO_K/\OO$ is a finite abelian group. The order of this quotient, denoted $[\OO_K : \OO]$, is 
called the \emph{index} of $\OO$ in $\OO_K$.
It is known that $m\OO_K\subset \OO$ for $m=[\OO_K : \OO]$.
For example, $\ZZ[2i]=\ZZ+\ZZ2i$ is an order of the Gaussian integers of index $4$, and $4\ZZ[i] \subset \ZZ[2i]$.


Let $p(x) \in K[x]$ be a polynomial. The \emph{splitting field} of $p(x)$ over~$K$ is the smallest extension of~$K$ over which $p(x)$ can be decomposed into linear factors. The  splitting field of $x^2-1$ over~$\QQ$ is $\QQ(i)$.

A \emph{root of unity} is any complex number that yields 1 when raised to some positive integer power $n$, \emph{i.e.}, $\zeta$ such that $\zeta^n = 1$. If $\zeta_n$ is an $n$-th root of unity and for each $k < n$, $\zeta^k \neq 1$, then we call it a \emph{primitive $n$-th root of unity}. We can always choose a primitive $n$-th root of unity by setting $\zeta_n = e^{2i\pi\frac{k}{n}}$ for $k$ with $k\in \ZZ_n^*$.
The $n$-th \emph{cyclotomic polynomial}, for any positive integer~$n$, is the unique irreducible polynomial $\Phi(x) \in \QQ[x]$ with integer coefficients that is a divisor of~$x^n - 1$ and is not a divisor of $x^k - 1$ for any $k < n$. The $n$-th cyclotomic polynomial~$\Phi_n$ is the minimal polynomial of a primitive $n$-th root of unity, and its roots are all $n$-th primitive roots of unity.
 
 \subsection{Galois theory}

An algebraic field extension $K/L$ is \emph{normal} (in other words, $K$ is normal over $L$) if every irreducible polynomial over $L$ that has at least one root in $K$ splits completely over $K$. In other words, if $\alpha\in K$, then all conjugates of $\alpha$ over $L$ (i.e., all roots of the minimal polynomial of $\alpha$ over $L$) belong to $K$.  An algebraic field extension $K/L$ is said to be a \emph{separable extension} if for every $\alpha\in K$, the minimal polynomial of $\alpha$ over $L$ is a separable polynomial. That is, it has no repeated roots in any extension field. Every algebraic extension of a field of characteristic 0 is normal. A \emph{Galois extension} is an algebraic field extension that is normal and separable. The following holds for separable extensions:

\begin{theorem}[Primitive Element Theorem]
\label{th:primitive_element}
Let $K/L$ a separable extension of finite degree. Then $K = L(\alpha)$ for some $\alpha \in K$; that is, the extension is simple and $\alpha$ is a primitive element. 
\end{theorem}
 
Given a Galois extension $K/L$, the \emph{Galois group} of $K/L$, denoted by $Gal(K/L)$ is the group of automorphisms of $K$ that fix $L$. That is, the group of all isomorphisms $\sigma: K \to K$ such that $\sigma(x) = x $ for all $x \in L$.
If $K$ is a field with subfield $L\subset K$, the \emph{Galois closure} of $K$ over $L$ is the field generated by images of embeddings $K \to K$ that are the identity map on $L$.

Fix $\alpha$ an algebraic number $\alpha$ over a Galois extension $K/L$.
The image of $\alpha$ under an automorphism $\sigma\in \Gal(K/L)$ is called a \emph{Galois conjugate} of $\alpha$. The Galois conjugates of $\alpha$  are precisely the roots of the minimal polynomial $f_\alpha$ of $\alpha$. 
The Galois conjugates of a root of unity~$\zeta_n$ are its powers~$\zeta^k_n$ such that $k\in \ZZ^{*}_n$; and $Gal(\QQ(\zeta_n)/\QQ)$ includes all automorphisms $\sigma$ defined by $\sigma(\zeta_n) = \zeta_n^k$ for  $k\in \ZZ^*_n$.



The \emph{norm} of $\alpha$ is defined by
\[ N_{K/L}(\alpha) = \prod_{\sigma \in \Gal(K/L)} \sigma(\alpha)\]
For short, we may drop the subscript $K/L$ if the underlying field is understood from the
context. 
For $\alpha = a + bi \in\ZZ[i]$ the only Galois conjugate is $a - bi$, and thus  its norm is the product
$N(\alpha) = (a + bi)(a - bi) = a^2 + b^2$. Note that the norms of all Galois conjugate are equal, and the norm of an algebraic integer is always a rational integer itself.

The \emph{trace} of $\alpha \in K/L$ is defined by:
\[ \Tr_{K/L}(\alpha) = \sum_{\sigma \in \Gal(K/L)} \sigma(\alpha) \]
Again, we drop the subscript $K/L$ if the underlying field can be understood from the context.

The ring of integers of $K$, $\OO_K$, is a free abelian group of rank $n$, and hence admits $\ZZ$-basis $\{\alpha_1, \ldots, \alpha_n\}$. Given such a basis, we denote with $\Delta_K$ the \emph{discriminant}, and define it by
\[ \Delta_K = \det(\Tr_{K/L}(\alpha_i\alpha_j))_{1\leq i,j \leq n}\]
Note that $\Delta_K$ is always a non-zero rational integer.





\subsection{Ramification theory}
Let $K$ be a number field and let $\mf{p}$ be a prime ideal of $\OO_K$. Then $\mf{p} \cap \ZZ$ is a prime ideal of $\ZZ$, hence there must exist a rational prime  $p$ such that $\mf{p} \cap \ZZ = p\ZZ$. We say that $\mf{p}$ \emph{is above} $p$. We have:
\begin{align*} 
\mf{p} \ \subset \  & \OO_K \  \subset \ K \\
 &~\mid \\
 p \ \subset &\ \ZZ \ \subset \ \QQ. 
\end{align*}

Given a number field $K$ with ring of integers $\OO_K$, any ideal $I\subseteq\OO_K$ admits a unique factorisation into prime ideals in $\OO_K$. Let $p\in\ZZ$ be a rational prime. The ideal~$p\OO_K$ may not be prime in $\OO_K$, but does factorise into prime ideals as follows:
\begin{equation}\label{eq:ram} p\OO_K = \mf{p}_1^{e_1} \cdots \mf{p}_g^{e_g}. 
\end{equation}


Using the vocabulary introduced above, we can observe that the prime ideals $\mf{p}_i$ are all above $p$. 
Note that in the ring of integers of a number field, all prime ideals are maximal, hence all $\mf{p}_i$ are also maximal ideals of $\OO_K$.
In general, given a commutative ring $R$ and a maximal ideal $\mf{m}$ of $R$, the \emph{residue field} is the quotient $k = R/\mf{m}$. Intuitively, $k$ can be thought of as the field of possible remainders. Now, given a maximal ideal $\mf{p}$ of $\OO_K$, $\OO_K/\mf{p}$ is a $\FF_p$-vector space of finite dimension. The \emph{residual class degree} (\emph{inertial degree}), denoted $f_\mf{p}$, is the dimension of the $\FF_p$-vector space $\OO_K/\mf{p}$, that is:
\[ f_\mf{p} = \dim_{\FF_p}(\OO_K/\mf{p}).\]
Looking at the factorisation \eqref{eq:ram}, we can define the residue class degree of each one of the $\mf{p}_i$'s as follows:
    \[ f_i = [ \OO_K/\mf{p}_i : \ZZ/p\ZZ ]. \]

We say that $p$ is \emph{ramified} if the \emph{ramification index} $e_i > 1$ for some $\mf{p}_i$. A prime $p$ is said to be \emph{totally ramified} if $e=n$, $g=1$, and $f=1$. That is $p\OO_K = \mf{p}^{e}$ for some $\mf{p}$.
 Conversely, $p$ is \emph{non-ramified} if $p\OO_K = \mf{p}_1 \cdots \mf{p}_g$ where the $\mf{p}_i$ are distinct. We further say that a prime $p\in\ZZ$ is \emph{inert} if the ideal $p\OO_K$ is prime, in which case we have $p\OO_K = \mf{p}$, that is $g=1$, $e=1$, and $f=n$.
 Finally, a prime is said to be \emph{split} if $e_1 = \ldots = e_g = 1$.
 
For the Gaussian integers, the ideals $2\ZZ[i]$ and $5\ZZ[i]$ are not  prime ideals and have respective factorisations $2\ZZ[i]=\mf{p}^2$ and $5\ZZ[i]=\mf{p}_1\mf{p}_2$ where $\mf{p}=(1+i)\ZZ[i]$, $\mf{p}_1=(2+i)\ZZ[i]$, and $\mf{p}_2=(2-i)\ZZ[i]$ are prime ideals. The prime $2$ is the unique ramified prime in the Gaussian integers.



\subsection{The $p$-adic field $\QQ_p$}

Here, we give a brief preliminary on the field of $p$-adic numbers $\QQ_p$; for more details see, e.g., \cite{gouvea2003p}. 
The field $\QQ_p$ is the completion of $\mathbb{Q}$ with respect to the $p$-adic absolute value $|\cdot|_p$, given by 
$|a/b|_p = p^{v_p(b)-v_p(a)}$ for $a,b\in \mathbb{Z}\setminus \{0\}$,
where $v_p(x)$ denotes the order to which $p$ divides $x \in \mathbb{Z}$. We denote by $\mathbb{Z}_p$ the valuation ring 
$\{ x \in \mathbb{Q}_p : |x|_p \leq 1\}$.
This is the local ring with unique maximal ideal generated by $p$.
A basic result about $\mathbb{Q}_p$ is Hensel's Lemma:

\begin{lemma}
\label{lem:hensel}[Hensel's lemma]
	Given $f(X) \in \mathbb{Z}[X]$, if there exists $\alpha \in \mathbb{F}_p$ such that 
	\[f(\alpha)=0 \text{ and } f'(\alpha)\neq 0\]
then there exists $x \in \mathbb{Z}_p$ with $f(x)=0$ and $x\equiv \alpha \bmod p$.
\end{lemma}

Given a number field $K/\QQ$, let $p$ be a rational prime and $\mathfrak{p}$ a prime ideal of $\OO_K$ lying above $p$. Then the $p$-adic absolute value $|\cdot|_p$ corresponding to $p$ extends uniquely to an absolute value $|\cdot|_\mathfrak{p}$ corresponding to $\mathfrak{p}$ such that the restriction of $|\cdot|_\mathfrak{p}$  to $\QQ$ coincides with $|\cdot|_p$.
This, in turn, corresponds to a field extension $K_{\mathfrak{p}}/\QQ_p$. This extension can be analysed using the ramification of $p$ in $K$. In particular, if $p$ completely splits in $K$, then $[K_{\mathfrak{p}} : \QQ_p] = 1$, that is, the extension is trivial and we have $K_{\mathfrak{p}} = \QQ_p$. If $p$ is inert, then the degree of the extension $K_{\mathfrak{p}}$ over $\QQ_p$ is equal to the inertial degree of $p$ in $K$. Finally, if $p$ is totally ramified, then the degree of the extension $K_{\mathfrak{p}}$ over $\QQ_p$ is equal to the degree of $K$ over $\QQ$, i.e., $[K_{\mathfrak{p}} : \QQ_p] = [K : \QQ]$.

Given a prime ideal $\mathfrak{p}$ of $\OO_K$, we define the decomposition group $D_{\mathfrak{p}}$ of to be the set of all automorphisms $\Gal(K/\QQ)$ fixing $\mathfrak{p}$, that is $ D_{\mathfrak{p}} = \{ \sigma \in \Gal(K/\QQ) \mid \sigma(\mathfrak{p}) = \mathfrak{p} \}$.
If the field $K$ is Galois over $\QQ$, the following isomorphism holds:
\[ D_{\mathfrak{p}} \cong \Gal(K_{\mathfrak{p}}/\QQ_p). \]
This entails that the prime $p$ completely splits in $K$ if and only if the decomposition group $D_{\mathfrak{p}}$ is trivial for all prime factors $\mathfrak{p}$ of $p$ in $\OO_K$.
